% =====================================================================
%  BACK MATTER
%  Conclusion, glossary, notes and references, further reading, index,
%  acknowledgments, about the author, colophon.
% =====================================================================

\backmatter

% ---- What's Next ----------------------------------------------------------
\chapter*{What's Next}
\addcontentsline{toc}{chapter}{What's Next}
\markboth{What's Next}{What's Next}
\thispagestyle{plain}

The most likely near future is not a resolution. It is a review.

The claim under examination has cleared the mechanical stages and stalled at
the mathematical one. That is the ordinary condition of a large formalisation,
and the machinery for the review exists --- the protocol of Chapter~\ref{ch:protocol}
is a summary of it --- but the work has to be done by people with the relevant
expertise, and it is months of reading per input.

Three specific things would move the situation, and they are worth watching
for.

\begin{enumerate}
\item \textbf{A written account of the mathematics.} The remark's author says
a prose write-up is routine and has not been written. When one appears, it
will resolve the 15-to-16 November gap in the only way that matters, and it
will make the review possible in the way that a 175-thousand-file development
does not.

\item \textbf{An independent statement review.} Someone in automorphic forms
and Iwasawa theory reading \texttt{Af.Hypotheses\_2}, \texttt{TLiftK} and
\texttt{admitted} against the published theorems. This is the step that
decides the question, and it is currently the step nobody has taken.

\item \textbf{An extension of the August result.} If the method can be pushed
past two thirds --- even to three quarters --- that would be independent
evidence about the ceiling, and it would be checked by people who are not
involved in the claim.
\end{enumerate}

What would not move the situation is more coverage. The story has been
covered extensively, and coverage cannot distinguish the three possibilities
in Chapter~\ref{ch:verdict}, because all three are consistent with everything
that has been said in public.

\vspace{10pt}
\noindent{\small\color{txtdim}\sffamily
For the reader who wants to follow this: the claim's own words are the best
primary source, the mirror's README states the verification standard precisely,
and the two-thirds paper is the best available evidence about the method. All
three are in the Notes and References.\par}

% ---- Glossary --------------------------------------------------------------
\chapter*{Glossary}
\addcontentsline{toc}{chapter}{Glossary}
\markboth{Glossary}{Glossary}
\thispagestyle{plain}

\noindent\small
Every term below is defined here in the sense in which the book uses it. Terms
marked \(\dagger\) are used in the sense given by the cited source rather than
in a general sense.

\vspace{8pt}
\begin{description}[style=nextline,leftmargin=0pt,font=\sffamily\bfseries\small,
  itemsep=5pt]

\item[Analytic continuation] An extension of a function to a larger domain,
agreed with the original wherever they overlap. For $\zeta$, the extension
from $\Re s > 1$ to all of $\CC$ is what makes the zeros in the critical strip
accessible at all.

\item[Automorphic form] A holomorphic function on an arithmetic group
satisfying a functional equation under that group, whose $L$-function is then
an $L$-function of the group. The theory of automorphic forms supplies the
$L$-functions whose zeros are studied in this book.

\item[Automorphy] The property of a representation of a Galois group of
arising from an automorphic form. Proving that a given Galois representation
is automorphic is one of the hardest known problems in number theory, and it
is one of the load-bearing steps of the method under examination.

\item[Axis] A vertical or horizontal line in the complex plane. The
\emph{critical axis} of $\zeta$ is $\Re s = \tfrac12$.

\item[CM field] A number field $E$ that is a quadratic extension $F(i)$ of a
totally real field $F$, with complex conjugation. CM fields carry rigid
analytic structures that make the divisibility statements of Chapter~5
possible.

\item[Compiler] The program that turns Lean source into kernel terms and
checks them. In the Lean 4 toolchain, \texttt{lean} performs elaboration and
the kernel performs verification.

\item[Critical line] The vertical line $\Re s = \tfrac12$ in the complex
plane. The Riemann hypothesis is the assertion that all the non-trivial zeros
of $\zeta$ lie on it.

\item[Critical strip] The band $0 < \Re s < 1$, containing all the non-trivial
zeros of $\zeta$.

\item[Dirichlet character] A homomorphism from $(\ZZ/N\ZZ)^\times$ to
$\mathbb{C}^\times$, giving the $L$-function $L(s,\chi) = \sum \chi(n)n^{-s}$.
The Riemann hypothesis is a statement about $\zeta$, which is the $L$-function
of the trivial character; the generalised version concerns all of them.

\item[Elaboration] Lean's process of turning human-written syntax into
kernel-checkable terms. Elaborated output is trusted for well-typedness, not
for correctness, which is why independent kernel replay exists.

\item[Euler--Maclaurin summation] The formula used to evaluate $\zeta(s)$ in
the half-plane $\Re s > 1$ by a finite sum plus an integral plus Bernoulli
corrections. The figures in this book use it.

\item[Functional equation] The identity $\Lambda(s) = \Lambda(1-s)$ relating
the completed zeta function $\Lambda(s) = \pi^{-s/2}\Gamma(s/2)\zeta(s)$ to
its reflection. It is the mechanism by which a non-vanishing result on one
side of the critical line is transferred to the other.

\item[Guinand's theorem] Refinements of the moments of $\zeta$ on the critical
line, by Édouard Guinand, published in the \emph{Proceedings of the London
Mathematical Society} in 1949 and 1950.

\item[Harmonic analysis on $\mathrm{GL}_2$] The theory, originating with
Hecke, of the automorphic forms and $L$-functions of $\mathrm{GL}_2$; the
source of the local factors $\det(1 + XA_y)/\det(1 - A_y)$ in the Lefschetz
formula of Chapter~7.

\item[Iwasawa algebra] A one-dimensional local algebra of the form
$\ZZ_p[[T]]$, occurring as the characteristic ideal of an Iwasawa module. The
Iwasawa tower of the submission supplies the arithmetic growth that the
integers otherwise lack.

\item[Iwasawa tower] A tower of number fields, or of curves over finite
fields, whose Iwasawa algebras are controlled. Behrend's construction of
such towers for function fields is invoked in the remark's item 19.

\item[Kato's Euler system] A class of Euler systems for modular forms
developed by Masaki Kato, whose existence is one of the inputs to the
admissibility and vanishing packages of Chapter~8.

\item[Kernel] The small, trusted part of a proof assistant that checks terms
against inference rules. Lean's kernel is roughly a thousand lines.

\item[Lefschetz fixed-point formula] The identity relating the alternating sum
of traces of maps induced on cohomology to the fixed points of the map. The
étale version used in the submission is proved in
\code{EtaleClopenDescent/LefschetzFixedPointFiniteEtaleCorrespondence.lean}.

\item[Level] A positive integer $N$ such that a modular form has good
reduction away from $N$. A condition on the level is a condition on the
arithmetic side.

\item[Mathlib] The large community library of mathematics formalised in Lean.
The submission is pinned to a specific revision of it.

\item[Multiplicity] The exponent of a point in a divisor. In the
Stepanov--Bombieri method, multiplicity is what converts a geometric dimension
count into an arithmetic bound; over $\ZZ$, it is the missing ingredient, and
supplying it is the central idea of the method under examination.

\item[Newform] A normalised cuspidal eigenform of weight $2$ at a level, whose
$L$-function has Euler product over the primes not dividing the level.

\item[\texttt{opaque}] A Lean declaration whose definition is not unfolded by
the kernel. Sound in principle; a place to look in an audit.

\item[Palomar registry] A registry of formalisations and their verification
status, intended to let claims about a formalisation be checked against a
third-party record.

\item[Poisson summation] The identity
$\sum_{n \in \ZZ} f(n) = \sum_{k \in \ZZ} \hat f(k)$ that underlies the
Parseval arguments of the remark's item 5.

\item[Proof assistant] A computer program that checks mathematical arguments
in a formal language.

\item[Rallis seesaw] An identity relating period integrals on one group to
period integrals on another, after which the period of the iterated theta
lift becomes a Bessel period independent of the Hecke characters.

\item[Riemann hypothesis] The assertion that every non-trivial zero of the
Riemann zeta function has real part exactly $\tfrac12$.

\item[Riemann zeta function] The analytic continuation of
$\zeta(s) = \sum_{n \ge 1} n^{-s}$ from the half-plane $\Re s > 1$ to the
whole complex plane, meromorphic with a single simple pole at $s = 1$.

\item[Riemann--von Mangoldt formula] The formula giving the number of zeros of
$\zeta$ up to height $T$:
$N(T) = \frac{T}{2\pi}\log\frac{T}{2\pi} - \frac{T}{2\pi} + O(\log T)$.

\item[Selmer group] A subgroup of a Galois module consisting of elements
satisfying local conditions at every place. Its vanishing is the content of
the statement \texttt{VanishOrd} in Chapter~8.

\item[Siegel--Weil formula] The formula relating the Fourier expansion of a
theta series to an $L$-function. It is one of the classical inputs to the
bundle \texttt{Tools\_k}.

\item[Statement review] The human process of checking that the formal
statements in a development mean what the accompanying prose claims they
mean. It is the step that a compiler does not perform, and it is the step
outstanding in this case.

\item[Structure constant] The number, in a finite-dimensional representation,
of the coefficient of the identity in a product. The GGP periods and
interlacing weights of the remark's item 7 are structure constants of this
kind.

\item[Submodule $\mathbb{C}^{\mathrm{UA}(E,n)}$] The space of unitary
automorphic representations of the unitary group $\mathrm{U}(E,n)$; the space
in which the theta lift of the submission takes place.

\item[Theta lift] An operation producing a new automorphic representation from
an old one, via a theta series attached to a Hermitian form. The iterated
theta lifts of the remark's item 8 are the method's central device.

\item[Trivial zeros] The zeros of $\zeta$ at $s = 0, -2, -4, \dots$, which
come from the poles of the gamma factor rather than from $\zeta$ itself, and
which the Riemann hypothesis excepts.

\item[Vandiver's conjecture] The conjecture that the plus part of the class
group of the maximal real subfield of a cyclotomic field has odd order. Open;
verified for primes below $2^{31}$ by Hart, Harvey and Ong in 2017.

\item[von Mangoldt function] The function $\Lambda(n)$, equal to $\log p$ if
$n$ is a power of a prime $p$ and $0$ otherwise. Its summatory function
$\psi(x) = \sum_{n \le x}\Lambda(n)$ encodes the primes.

\item[Weil's criterion] The positivity condition
$\sum_\rho \hat g(\rho)\hat g(1-\rho) \ge 0$, for all admissible test
functions $g$, which is equivalent to the Riemann hypothesis. Derived from
Weil's 1952 paper.

\item[Weil expression] The real-valued function $W(t)$ built from the
ordinates of the zeros of $\zeta$, whose non-negativity is equivalent to the
Riemann hypothesis, and whose positivity the remark seeks to obtain via a
determinant.

\item[$\Theta$] The supremum of the real parts of the zeros of $\zeta$.
Landau's oscillation theorem gives $\psi(x) - x = \Omega_\pm(x^{\Theta-\varepsilon})$,
and RH is the assertion $\Theta = \tfrac12$.

\end{description}

% ---- Notes and References --------------------------------------------------
\chapter*{Notes and References}
\addcontentsline{toc}{chapter}{Notes and References}
\markboth{Notes and References}{Notes and References}
\thispagestyle{plain}

\noindent\small
Full citations for every source relied on in the text. Where a claim could
not be tied to a primary record, the omission is itself noted. Entries marked
\emph{[artifact]} refer to files in the submitted Lean development, quoted by
their path relative to the archive root.

\vspace{10pt}
{\sffamily\bfseries\small Primary artifact}

\begin{enumerate}[leftmargin=1.4em,itemsep=4pt]
\item \textbf{Alpöge, Levent.} \emph{Remark}. Four pages, unpublished;
distributed at \texttt{alpo.ge/remark.pdf} and mirrored on 20 November 2026.
[artifact] Quoted throughout Parts~I and~II.
\item \textbf{Submitted Lean development.} \code{alpo.ge/lean.tar.gz.part01}
through \texttt{.part14}; 239\,275\,554 bytes; mirrored as
\code{github.com/CocoSgt/riemann-hypothesis-lean}, release \texttt{v2026-11-20},
SHA-256 \code{9631b4f020199e580ba86178cb56c6e9b09f690c2369a80dc4334aa93f49c7fc}.
[artifact] All Lean quotations in Chapters~7--9.
\item \textbf{README of the mirror repository.} \emph{Status}: ``As of mirroring
(2026-11-20), no independent build or review of this code has been published.
Treat it as a claimed proof until it has been compiled and checked.'' [artifact]
Quoted in Chapters~1 and~2 with its date, which matters.
\item \textbf{\texttt{lean-toolchain}} [artifact]: \code{leanprover/lean4:v4.33.1}.
\item \textbf{\texttt{lakefile.toml}} [artifact]: Mathlib pinned at
\code{0df444a360eaa60ab8c11dca51a86af692955474}; confirmed locally.
\item \textbf{\texttt{src/Main.lean}}, \texttt{src/Statements.lean} [artifact]:
quoted in full in Chapter~2.
\item \textbf{\code{CharacterBoxes/MainConclusion.lean}},
\code{CharacterBoxes/EveryCharacter.lean},
\code{CharacterBoxes/RiemannHypothesisOfEveryCharacter.lean} [artifact]:
quoted in Chapter~7.
\item \textbf{\code{Foundations/GaloisRepresentations/IntegralStructures/RowSrc.lean}},
\code{GaloisRepresentationsOfThetaLift.lean},
\texttt{Development.lean},
\code{GeometricFlagBases/WeightOneGeometricRealization.lean} [artifact]:
quoted in Chapter~7.
\item \textbf{\code{AbelianVarietyTheory/Proofs/SatakeParameters/LefschetzFixedPointFiniteEtaleCorrespondence.lean}}
and \code{AutomorphicArithmetic/Foundations/EtaleClopenDescent/LefschetzFixedPointFiniteEtaleCorrespondence.lean}
[artifact]: Chapter~7.
\item \textbf{\code{.../ThetaParseval/HalfIndexLattices/TLiftK.lean}},
\code{.../ThetaKernels/TLiftK.lean}, \code{.../ThetaKernels/TPGHolds.lean},
\code{.../ThetaKernels/ThkEstEpsOf.lean} [artifact]: Chapter~8.
\item \textbf{\code{AutomorphicVarieties/AutomorphicForms/Hypotheses2.lean}},
\code{.../TemperedLocalComponents/NewformAut.lean},
\code{.../ModularityArgument/ClozelCohomology/ClozelConjC.lean},
\code{.../PublishedTheorems/GelfandKazhdanReal.lean} [artifact]: Chapter~8.
\item \textbf{\code{AbelianVarietyTheory/Definitions/Foundations/FinalTheorems.lean}}
[artifact]: the structure \code{GlobalTheory.Ord2.admitted} and the statement
\texttt{VanishOrd}; Chapter~8.
\item \textbf{\code{AbelianVarietyTheory/Proofs/GlobalInterpolation/NonemptyHypsOfAllRows.lean}}
[artifact]: the construction of the \texttt{Hyps} package; Chapter~8.
\end{enumerate}

\vspace{8pt}
{\sffamily\bfseries\small The mathematics}

\begin{enumerate}[leftmargin=1.4em,itemsep=4pt,start=14]
\item \textbf{Alpöge, Levent, and Ralph Furman.} \emph{More than two thirds of
the zeta zeros are simple and on the critical line}. arXiv:2608.13637,
13 August 2026. Preprint, 33 pp. with appendices. Accompanying Lean
formalisation at \code{github.com/anthropics/formal-math} (formerly
\code{anthropics/zeta-23-lean}).
\item \textbf{Lamzouri, Yousouf.} arXiv:2609.02882, 2026. Independently
reproduces the constants of the preceding paper by a different route and
credits ``an internal research version of Claude developed by Anthropic'' and
``two mathematicians at Anthropic, Levent Alpöge and Ralph Furman''. The
precise title was not independently confirmed in this book's research and is
not asserted here.
\item \textbf{van Frankenhuijsen, Maarten.} \emph{The Riemann Hypothesis for
Function Fields over a Finite Field}. arXiv:0806.0044, 2008. The source of the
statement that the Stepanov and Bombieri methods require a product $C \times C$
with one geometric and one arithmetic coordinate, and that over $\ZZ$ the
product $\operatorname{Spec}\ZZ \times \operatorname{Spec}\ZZ =
\operatorname{Spec}\ZZ$ is one-dimensional.
\item \textbf{Weil, André.} The criterion for the Riemann hypothesis as a
positivity condition on a quadratic form, in the paper of 1952 published from
Lund. Confirmed via Burnol, arXiv:math/9810169, and Groskin, arXiv:2607.02828.
\item \textbf{Guinand, Édouard.} \emph{Proceedings of the London Mathematical
Society}, (3) \textbf{1} (1949) and (3) \textbf{2} (1950). The frequently
repeated 1955 title \emph{Sur les zéros de la fonction zêta de Riemann} could
not be verified and is not cited here.
\item \textbf{Titchmarsh, Edward C.} \emph{The Theory of the Riemann
Zeta-function}, 2nd ed., revised by D. R. Heath-Brown. Clarendon Press, 1986.
Theorem 15.11 for the quantitative oscillation $\Omega_\pm(x^{1/2}\log\log\log x)$.
\item \textbf{Landau, Edmund.} The oscillation theorem for $\psi(x)$, in the
general Dirichlet-integral form; see Titchmarsh, above, and
Iwaniec--Kowalski, \emph{Analytic Number Theory}.
\item \textbf{Bost, Albert H.} \emph{La conjecture de Weil: preuve d'un théorème
de D. Bertrand et P. Swinnerton-Dyer}. \emph{Annales scientifiques de l'École
Normale Supérieure} \textbf{30} (1997), no.~1, 229--288.
DOI \code{10.1016/S0012-9593(97)89923-8}. The precise ``slope inequality''
as used in the argument was not retrieved from a primary text and is therefore
paraphrased rather than quoted in this book.
\item \textbf{Behrend, Theodor.} The construction of Iwasawa towers for
function fields, 1940. No DOI or OpenAlex record was retrievable in this
book's research, and no identifier is asserted. The number-field analogue
remains open for $d \ge 2$.
\item \textbf{Hart, G. L., D. Harvey, and L. Ong.} \emph{Irregular primes to two
billion}. \emph{Mathematics of Computation} \textbf{86} (2017), 3031--3049.
DOI \texttt{10.1090/mcom/3211}; arXiv:1605.02398. The bound $2^{31}$ on
Vandiver's conjecture.
\item \textbf{Gan, Wee-Tak, and Ichino, Ai.} The unitary Gross--Prasad
theorem, and the related results for theta lifts. \emph{The Gross--Prasad
theorem and local theta correspondence}; the orthogonal case remains open.
The commonly used attribution ``Gross--Petersson'' is a misattribution and is
not used here.
\item \textbf{Jacquet--Shalika rigidity}, \textbf{Borel injectivity}, the
\textbf{Gelfand--Kazhdan correspondence} for $\mathrm{GL}_2$, and
\textbf{Kato's Euler systems} are the published results underlying the
\texttt{Hypotheses\_2} bundle of Chapter~8. Individual primary citations for
each were not separately retrieved for this book, and the bundle's own
derivation chains in the artifact are the reference given in the text.
\end{enumerate}

\vspace{8pt}
{\sffamily\bfseries\small The mathematics of the problem}

\begin{enumerate}[leftmargin=1.4em,itemsep=4pt,start=26]
\item \textbf{Mathlib.} \code{Mathlib/NumberTheory/LSeries/RiemannZeta.lean}
at revision \code{0df444a360eaa60ab8c11dca51a86af692955474}:
\texttt{riemannZeta} at line 121, defined as \texttt{hurwitzZetaEven 0}, the
analytic continuation rather than the Euler product; \texttt{RiemannHypothesis}
at lines 185--186. The Euler-product zeta function is separate, as
\code{zeta : ArithmeticFunction} $\mathbb{N}$. Both definitions were read directly
from the source at the pinned commit.
\item \textbf{The Riemann explicit formula}, the equivalence of RH with
$\psi(x) = x + O(x^{1/2}\log^2 x)$, and the equivalence with Weil's criterion:
standard; see Titchmarsh and Iwaniec--Kowalski, above.
\item \textbf{Pratt, Robby, Matthew Robles, and Kamal Zaharescu.} Bounds on the
proportion of zeros on the critical line: 41.7293 per cent on the line and
40.7511 per cent simple and on the line, as the 2020 records.
\item \textbf{Conrey, John B.} (1989) on the proportion of zeros on the
critical line, as the earlier asymptotic record.
\end{enumerate}

\vspace{8pt}
{\sffamily\bfseries\small Formal verification, and the claim}

\begin{enumerate}[leftmargin=1.4em,itemsep=4pt,start=30]
\item \textbf{arXiv:2606.29493.} The explicit statement that the Lean kernel
verifies that a formal statement has a proof, and does not verify that the
formal statement matches the informal claim it is meant to capture.
\item \textbf{Lean FRO.} The Formalisation and Reasoning Overview, at
\code{lean-lang.org/fro/about/}. Documents the tools --- Comparator,
\texttt{lean4export}, \texttt{landrun}, kernel replay, Palomar, and the
\texttt{formalization.yaml} metadata convention. Cited in the arXiv
literature as ``Lean FRO, 2025/2026''. No peer-reviewed paper titled
``LeanFRO'' was located, and none is cited.
\item \textbf{\texttt{lean4checker}} (\code{github.com/leanprover/lean4checker}),
for third-party checking of a compiled environment.
\item \textbf{Quanta Magazine.} Kakaes, Konstantin. ``A claimed proof of the
Riemann hypothesis'', 20 November 2026. Source for the independent
compilations by Andrew Sutherland (MIT) and Simon DeDeo (CMU), and for the
caveat about Lean's being a programming language that ``can also be used to
play sneaky tricks by, for example, redefining symbols used in a proof to
make something trivially true''.
\item \textbf{Agence France-Presse.} Report of 20 November 2026: ``The claim
has not been independently verified, and the proof has not yet been reviewed by
outside experts.''
\item \textbf{OpenAI project history.} The record of a formalised result
withdrawn in October 2026 with a stated sign error --- cited as the precedent
for the distinction between compiling and being correct.
\end{enumerate}

\vspace{8pt}
{\sffamily\bfseries\small Deliberately not cited}

\begin{enumerate}[leftmargin=1.4em,itemsep=4pt,start=36]
\item \textbf{The ``Saouter--Demichel--Planchat'' paper and the ``LP method
with crossing-zero value $L = 0.40$''.} No such paper and no such value were
located in arXiv, Crossref, HAL or OpenAlex during this book's research. The
two-thirds result is attributed to arXiv:2608.13637, and no crossing value is
asserted anywhere in this book.
\item \textbf{Community discussion threads}, including Hacker News threads,
sceptical blog responses, and a separate announced GRH claim by another
party. These are evidence that discussion exists, and are not used as
mathematical sources or as evidence about the state of the mathematics.
\item \textbf{Secondary descriptions of the argument.} No account of the
mathematics has been taken from press coverage. Every description of the
method in this book derives either from the remark itself or from the
artifact's source.
\end{enumerate}

% ---- Further Reading -------------------------------------------------------
\chapter*{Further Reading}
\addcontentsline{toc}{chapter}{Further Reading}
\markboth{Further Reading}{Further Reading}
\thispagestyle{plain}

\noindent\small

\textbf{For the hypothesis itself.} Edward C. Titchmarsh, \emph{The Theory of
the Riemann Zeta-function}, 2nd ed., revised by D. R. Heath-Brown, Clarendon
Press, 1986, remains the standard single reference. For a shorter and more
reflective account: A. Ivić, \emph{The Riemann Zeta-Function}, and the
collected essays in \emph{The Riemann Zeta-Function and Its Derivatives}.

\textbf{For the function-field proof.} Maarten van Frankenhuijsen's
arXiv:0806.0044 is the most accessible account of why the Stepanov and
Bombieri methods work over finite fields and why the arithmetic analogue
fails. Weil's original proof is in \emph{Number Theory: An Introduction to
Mathematical Analysis of the Theory of Numbers} (1952).

\textbf{For the automorphic machinery.} The reference here is
\textsc{Automorphic Forms, Representations and $L$-functions}, Proceedings of
the 1989 Princeton conference, for the general setting; and the work of
Gan--Ichino, and of Liu and W. Zhang, for the unitary theta lifts and the
Gross--Prasad theorem in the unitary case. The distinction between the
proved unitary case and the open orthogonal case is the point of Chapter~4's
second barrier.

\textbf{For Iwasawa theory.} Coates--Su, \emph{Iwasawa Theory for Number
Fields}, for the general framework; Greenberg and Wan's work for the
function-field and number-field towers. Behrend's 1940 construction is the
origin of the idea invoked in the remark's item 19.

\textbf{For formalisation.} The Lean documentation; the Lean FRO pages at
\code{lean-lang.org/fro/about/}; and the Palomar registry, whose entries
record the verification status of formalisations including the two-thirds
theorem's. For the standard of what a formalised proof should be accompanied
by, arXiv:2606.29493 is the most directly relevant recent statement.

\textbf{For the claim under examination.} The remark itself, the mirror
repository's README, and the two-thirds paper. Everything else that has been
written about this claim is either a summary of those or a summary of the
press coverage, and this book has tried to use only the first kind.

% ---- Index -----------------------------------------------------------------
\printindex

% ---- Acknowledgments -------------------------------------------------------
\chapter*{Acknowledgments}
\addcontentsline{toc}{chapter}{Acknowledgments}
\markboth{Acknowledgments}{Acknowledgments}
\thispagestyle{plain}

This book exists because a remarkable document was made available in public,
and because the question it raises is more interesting than the answer anyone
currently has to it.

The mathematics in Chapter~3 and Chapter~4 is standard, and the people who
built it over two centuries --- Riemann, Hadamard, de la Vallée Poussin,
Hardy, Littlewood, Weil, Stepanov, Bombieri, Guinand, Bost, Behrend --- are
owed more than a footnote each. Their results are the terrain; the claim under
examination is an attempt at a route across it, and the route is only
intelligible because of the terrain.

The claim itself belongs to Levent Alpöge, who made the artifact available
under his own name and who wrote a document that is, to this book's lasting
impression, more candid about its own limits than most published mathematics.
The four-page remark deserves to be read as evidence rather than as a
news hook, and the passages where it says ``the Lean should suffice to
clarify everything'' and ``I'm just too inexpert in the relevant stuff at the
moment'' are the reason this book could be written as an audit rather than as
an obituary for a claim.

Thanks are due to the compilers --- Andrew Sutherland and Simon DeDeo --- who
did the unglamorous work of finding out whether the artifact builds, and to
the mirror's maintainer, whose README is a better piece of technical writing
than most published formalisation papers.

The verification work in Chapters~7 to~9 was done on the submitted source
tree, and the corrections log that accompanies this book records the places
where the writer's initial assumptions had to be revised in the course of it.
Several of those revisions made the book weaker in tone and stronger in fact,
which is the only trade this book regards as a good one.

\begin{flushright}
\emph{Gerald Enrique Nelson Mc Kenzie}\\
Tico Internet LLC\\
21 November 2026
\end{flushright}

% ---- About the author ------------------------------------------------------
\chapter*{About the Author}
\addcontentsline{toc}{chapter}{About the Author}
\markboth{About the Author}{About the Author}
\thispagestyle{plain}

\noindent\textbf{Gerald Enrique Nelson Mc Kenzie}\par
\vspace{6pt}

\noindent\small
Gerald Enrique Nelson Mc Kenzie is the author of \emph{Uniworld: A Relational
Monism of Inseparable Frames} (Tico Internet LLC, 2026 First Edition, Apache
License 2.0, DOI 10.5281/zenodo.24119718). He writes at the intersection of
mathematics, formal verification and the assessment of machine-generated
claims, and he is not a number theorist, which this book says wherever it
matters.

His interest in this subject is practical rather than decorative. A proof
claim that arrives as seventeen million lines of code is a new kind of
evidentiary object, and the question of how to read one --- what to check
first, what a compiler settles, what it cannot settle, and in what order ---
is a question that will recur for every important formalisation from now on.
The answer he has arrived at is that the mathematics has to be read, and that
saying so is more useful than either applause or dismissal.

He publishes under the imprint Tico Internet LLC.

\vfill
\noindent{\sffamily\small\color{txtdim}\rule{\textwidth}{0.4pt}}\par
\noindent{\sffamily\footnotesize\color{txtdim}
\emph{Corrections.} Corrections to this edition are welcomed and will be
carried in a later printing. The corrections log accompanying the research
materials for this title records, in the open, the places where the writer's
initial assumptions were revised during preparation.\par}

% ---- Colophon -------------------------------------------------------------
\clearpage
\thispagestyle{empty}
\vspace*{1.0in}
\begin{center}
\begin{minipage}{0.92\textwidth}
{\sffamily\fontsize{15}{18}\selectfont\bfseries\color{txt} \bookTitle\par}
\vspace{4pt}
{\sffamily\small\color{txtdim} \bookSubtitle\par}
\vspace{18pt}
{\color{accent}\rule{\textwidth}{1.6pt}}\par
\vspace{14pt}
{\small\color{txt}
\noindent\textbf{Author}\quad \bookAuthor\par
\vspace{4pt}
\noindent\textbf{Publisher}\quad \bookImprint\par
\vspace{4pt}
\noindent\textbf{Edition}\quad \bookEdition\par
\vspace{4pt}
\noindent\textbf{DOI}\quad \href{https://doi.org/\bookDOI}{\texttt{\bookDOI}}\par
\vspace{4pt}
\noindent\textbf{Licence}\quad \bookLicense\par
\vspace{4pt}
\noindent\textbf{Source code}\quad None published for this edition. The Lean
development examined in this book is the work of its author and is not
redistributed here.\par
\vspace{14pt}
\noindent\textbf{Subject headings}\quad
Riemann hypothesis --- Mathematical verification; Number theory ---
Computational methods; Formalization (Computer science); Machine learning ---
Mathematical applications; Proof assistants; Theta functions; Iwasawa theory.\par
\vspace{14pt}
\noindent\textbf{Errata}\quad
No errata have been recorded for this edition.\par
\vspace{14pt}
\noindent\textbf{Extent}\quad
\bookPages{} numbered pages, with full-bleed front and back
covers;\ \bookWords{} words.\par
\vspace{14pt}
\noindent\rule{\textwidth}{0.4pt}\par
\vspace{10pt}
\noindent\textbf{Signed}\quad
\emph{Gerald Enrique Nelson Mc Kenzie}\par
\vspace{4pt}
\noindent Tico Internet LLC\par
\vspace{4pt}
\noindent \bookFactsDate\par
}
\end{minipage}
\end{center}
